
% 28 Mar: Clean up


% Sverdrup and the Necessity of the Local Congregation [editorial title]

\subsection{Installment 6: “Folkebladet,” April 17, 1918}

Whoever has read the significant work of the Swedish professor Gustaf Aulen from 1912, “The Lutheran Church-Idea,” in which the author deals with the Lutheran concept of the church as well as the ‘Lundensian’ and ‘Uppsalian’ concepts of the church, will have noticed a quotation in the preface from one of Scotland’s leading churchmen, John Cairns (died 1892). Cairns says:

\begin{quote}
With all that the labour of the theologians has done
to define and vindicate the idea of the church, I believe
that we are as yet only on the surface of its depths.
\end{quote}

I think Cairns would have appreciated Sohm’s treatment of the church. Yes, I think he would also have rejoiced over the “guiding principles,” though they do not go so far as Sohm.

One thing must now be clear. The guiding principles and that which their authors have written concerning them must be studied, discussed, and clarified—at least in private. If one immerses oneself in these principles, then a hundred questions will arise, and one will scarcely give up before one has received satisfactory answers from the authors. If, in addition, one reads related literature written by others, then new points of view emerge; and one then returns to the principles with these, makes comparisons, and forms for oneself a comprehensive view in which no details have been skipped over, nothing foreign has been imported, and the greater and the lesser have received their proper importance and their proper place.

It is not for everyone to do this work. Where prejudice has anchored itself fast, where intellectual alertness is dull, indifference great, and the assertion “We have Abraham as our father” strong, there the prospects for a fruitful engagement with the principles and the literature on church polity are by no means promising. But where there is good will, an honest endeavor to get hold of what is right, there one may come a long way. And it is needed.

Principle no. 1 especially needs our attention. Do we all understand it? Is it too much to ask whether it is fully in accord with the Word? Is it absolutely excluded that the authors could have been mistaken? The opposition to this principle at any rate justifies the raising of such questions. Thus: questions awaiting answers, tasks awaiting solution.

One may, of course, get away from these questions and tasks by once more asserting that principle no. 1 aims at the unorganized holy catholic Church, the communion of saints. But that is precisely what it does not do.

One may indeed make objections by pointing to the exposition of principle 4, which speaks of the congregation’s outward organization. But this principle must be understood in connection with principle no. 1.

Principle no. 1 is directed against two things. First, against the complex organizations which we call the organized church, the state churches, the denominations, and so forth. The principles regard these organizations not as church or congregation in the sense of Scripture, but as worldly organizations. Second, they are directed against the view that the communion of saints, the unorganized church, is all that is needed. For the Church in the sense of Scripture is presented by S. and D. as consisting of congregations, not individuals.

Already in 1875 S. and D. turned “the weapons ... against pietism and false spirituality, which dissolves the Christian community into isolated individuals.” “It belongs inseparably to the essence of the Christian Church to exist in the form of congregations, not as individual Christians.”

In 1895 S. and D. wrote (joint declaration): “The Church is the Kingdom of God on earth, which between Christ’s ascension and His coming in glory has no other rightful form than the congregation.”

In 1898 S. declared: “This is something new when in point 1 we declare the congregation to be the proper form of the Kingdom of God on earth; for for hundreds of years men have called the church the proper form.”

That this means something other than adopting a new name for an old thing appears from what S. wrote in 1896:

\begin{quote}
The first thing, then, which comes forward as the scriptural principle in the relation between congregation and church body is this: that the congregation is the original, and the church body the derived; the congregation is the true, proper, apostolic, and Christian form of the Kingdom of God on earth in the time from the outpouring of the Spirit until Christ’s coming again; all outward societies or churches or institutions or associations or companies have, according to God’s Word, their justification only as serving instruments for the congregation, not as authorities that govern and issue commands.

This principle is the first and most far-reaching of all rules for church polity. The individual local congregation has God’s Word over it, has Jesus Christ as head and Lord, has God’s Spirit as its guide; but it has no other authority over it.

It is entirely clear that the New Testament teaches that Christians should not only have an inward fellowship of the Spirit with one another, but that they should also form outwardly organized congregations with their officebearers and their rules. We should depart from the altogether clear demand of God’s Word if we sought to be so spiritual that we would have no outward congregation. The outward congregation belongs to apostolic and biblical Christianity. Among the many who refuse to belong to any outward congregation because it is not spiritual enough for them, some suffer from a grievous lack and weakness in their Christianity, though it is far from us to mean that such frail Christians are not Christians at all. On the other hand, the same does not hold for what is called church bodies ...

“The very first principle of the critique of church polity is therefore this, that the Christian Church consists of free, independent, and self-governing congregations ... Such is the teaching of Scripture. (Veiledning 335).

\end{quote}

The congregational form is necessary, according to S. “The Kingdom of God has its inner essence, which is righteousness, peace, and joy in the Holy Spirit. It also has its outward form, which is the congregation, the Spirit’s people in the world.” The congregation is this outward form, “the only outward form appointed by God for His Kingdom on earth.”

But this outward form of the Kingdom of God, the congregation, itself has both an outward side and an inward side.

There is talk of “its inward side in its invisible essence,” its “inner spiritual essence.” According to this, in “its innermost essence it [the congregation] is the body of Jesus Christ, or the work of the Holy Spirit.” “In the dogmatic sense, people are not received into the congregation by a decision of the congregational meeting—but by the Word, Baptism, and faith.” (S. and D. 1895).

But the congregation also has its outward side, its mixed character, the formal organization.

It is the believers in the local congregation who are the body of Christ. The unbelieving members are members of the congregation in an improper sense; they stand only in an outward connection.

The proper form of the Kingdom is the outward congregation: “The individual outward congregation is the body of Christ ... every member in this outward congregation is a member of Christ’s body” (Oftedal in Veiledning 65). “The body of Christ ... is visible, because it is believing men and women who work with the Word and the Sacraments, and it is the outward, organized fellowship of these which is called congregation.” (S. in Veiledning, 38). “What is said about the unity of all congregations is also said about the individual outward congregation.” (Oftedal in Veiledning, 65).

And again Sverdrup:

\begin{quote}
For all congregations and for every one of them Jesus Christ is the living foundation upon which the congregation is built ...

“Paul did not mean that Jesus and His life, death, and resurrection are once for all the foundation of the Church, and since it has once been laid, it cannot be laid again; rather, he means that by the preaching of the Gospel in Corinth this foundation was laid there again, and the Corinthian congregation was built upon it.

“And not that alone. All congregations, and every single one of them, are the body of Jesus Christ. In this body He is the soul, and for this body He is the head ... As intimately as the soul is joined to the body, so near is Jesus to the congregation ... He Himself is the head and governs and leads the congregation, works through it. (Luth. Kirkeblad 1893, no. 4).

... every congregation and every congregational member must remember that it is not merely the whole great Church, the gathering of all congregations, which is the body of Jesus and the bride of Jesus; rather, it is precisely the little gathering of people who in each place gather around the Word and the Sacraments that has the calling to be God’s temple, the body of the risen and glorified Savior, through which He Himself will work on earth, by which He still reveals Himself on earth.

Therefore the responsibility is so great. God’s temple is holy; therefore it is all the greater a transgression for a congregational member to be profane and thereby corrupt God’s temple.

Luther’s words about ‘the whole Christian congregation on earth’ must also be true of each individual congregation, if this truly answers to its name. (L. K. 1893, no. 10).









\end{quote}

Some of the most unfortunate expressions theologians have been accustomed to use concerning the congregation are: visible and invisible, inward and outward. They are thoroughly confusing. S. and O. have to some degree accepted this usage, and it was scarcely to be expected that they could use these expressions for 35 years without committing one inconsistency or another—for who has not been made to stumble by these expressions? \footnote{In 1877 Oftedal wrote: The congregation is visible because God’s Word and the Sacraments are there, and invisible because the believers are there and it is “the body of Christ” (Kvartalskrift 1877, 15).}

One should note the difference.

There may be an outward congregation even where external organization is lacking. “It is the use of the Word and the Sacraments that constitutes a congregation in the biblical sense” \footnote{In 1897 Sverdrup wrote: The congregation becomes visible by its believing members and their activity. ... The body of Christ ... is visible. (Veiledning 37, 38).}. This means the local congregation—not “the whole Christian Church.” The outward organization, the outward congregational form, follows as a rule afterward.

\begin{quote}
Where these means [Word and Sacraments] are rightly used, there is at least a mission, and where these means work eternal life in souls, there the body of Christ and His congregation are present according to their essence, even though there is not yet any congregational organization” (S. S. II, 236).

“The congregation [the outward congregation] was there ... already at Pentecost. It had adopted no constitution, chosen no officebearers, organized itself in any way by outward regulations and decisions; nevertheless it was there.” (S. S. II, 9).

\end{quote}

The formal organization is the external, outward organization. But the organization in reality is the congregation, the local congregation.

(N. B. We distinguish sharply between outward organization and outward congregation.)

The local congregation, when it answers to its name, is a gathering of believers. It is defined in principle no. 2, which deals precisely with the congregation’s constitutive members. This principle states what the congregation theoretically is. Those who are called constitutive members are believers. But in the outward organization there are also hypocrites; they intrude into it, they are no constituent part, their connection is only outward. That is what the congregation in its practice is according to principle no. 4, and the Augsburg Confession, art. 8. (cf. S. S. III, 368).

S. and O. laid such weight on the constituted local congregation (organization may follow later) that some may say that they made it into a sacrament. And it is true that S.’s view of the congregation resembled very much the early church’s view of the sacraments. Therefore S. could say that one who lived alone on an island had no need of the Supper, the congregation’s meal. The older theologians said that it is not the lack of the Sacraments that condemns, but contempt for them. Compare Sverdrup: “Where an individual Christian lives alone in prison or under some other compulsion, there the lack of a congregation will not lead to his perdition (Veiledning 95).” “The congregation is absolutely necessary where it is possible.” “Where Christians can be congregation and will not, there is at times great danger for their soul (95).”

“For our Free Church people, the matter stands thus: that when God’s Word has set the congregation as the proper form for the Kingdom of God on earth until Christ’s coming, then it is both necessary and possible to have congregation, whatever it may cost” (against the free-free).

All this, of course, points to the “outward congregation.”

The local congregation was of such significance for S. and O. that, for example, in the 1870s they did not wish any local congregation to be established at Augsburg, although their colleagues Gunnersen and Gjeldaker wished this. The matter was dealt with several times in the faculty. Even a faculty opinion was submitted. Among the grounds S. and O. advanced was this, that school discipline with its strictness would interfere disruptively with the congregation as a community of professors (and their families) and students. Church discipline and school discipline might conflict with one another.

The local congregation was of such significance for S. and O. that they never stated, so far as I know, that an individual could be a member of several congregations at once; whereas they saw nothing wrong in a congregation’s standing in connection with several church bodies, since these (in contrast to the congregation) are not governed by the law of the body of Christ. The church bodies were only organs, worldly arrangements, which could serve the congregation, which was the actual organization, the divine institution, the only form in which the holy catholic church becomes manifest (cf. Pontoppidan F. C. B. edition, col. 243).

No matter how much one wishes to get away from the fact that S. and O. held this doctrine of the congregation; it is the same how much one wishes to twist their statements—one cannot get away from the concept of the local congregation as the fundamental thing in S. and O’s conception of the Church. \textbf{For them the local congregation was not merely something expedient, not merely an organization necessary in an organic sense; for them it was at least just as necessary in the specifically Christian sphere as the family is in the social. One may call the whole human race a family, but the individual family is a necessary ordinance for the proper growth and upbringing of the human race.}  Neither state nor school nor orphanage can wholly replace upbringing in the home. So too, all Christians together are called congregation, or church, but this whole Christian congregation or church consists of individual congregations.



